\documentclass[a4paper]{report}
% Shared preamble for the English probability and statistics documents.
% Based on the original course preamble. Compile with pdfLaTeX.
\usepackage[top=1.5cm,bottom=2cm,left=2cm,right=2cm]{geometry}
\usepackage[T1]{fontenc}
\usepackage{lmodern}
\usepackage[utf8]{inputenc}
\usepackage[english]{babel}

\usepackage{amsmath}
\usepackage{amssymb}
\usepackage{amsthm}
\usepackage{mathrsfs}
\usepackage{mathtools}
\usepackage{etoolbox}
\usepackage{tikz}
\usetikzlibrary{shapes.multipart,arrows.meta}
\usepackage{xcolor}
\usepackage{enumerate}
\usepackage[shortlabels]{enumitem}
\usepackage{hyperref}
\usepackage{version}
\usepackage{bbold}
\usepackage{fontawesome5}
\usepackage{fancyhdr}
\lfoot{}
\usepackage{natbib}
\usepackage{graphicx}
\usepackage{pdfpages}
\usepackage{authblk}
    
\usepackage{blindtext}
\usepackage{multicol}
\usepackage{collcell}
\usepackage{wrapfig}
\usepackage{blkarray}
\usepackage{cancel}
\usepackage{appendix}
\usepackage{float}
\makeatletter
\input{tex/assets/slashbox.sty}
\makeatother
\hypersetup{unicode,colorlinks,citecolor=red,filecolor=black,
  linkcolor=red,urlcolor=cyan}
\setlength{\columnsep}{1cm}
\pagenumbering{gobble}

\def\numtext#1{%
  \ifcase#1\or one\or two\or three\or four\or five\or six\or seven\or eight\or nine\or ten\or
  eleven\or twelve\or thirteen\or fourteen\or fifteen\or sixteen\or seventeen\or eighteen\or nineteen\or twenty\or Lots\fi}

\newcommand{\Byte}[1]{%
  \def\nodecontents{}%
  \foreach[count=\l] \k in {#1} {\xappto{\nodecontents}{\noexpand\nodepart{\numtext{\l}}\k}}%
  \node[rectangle split,rectangle split parts=8,rectangle split horizontal,draw] at (2,2){\nodecontents};}

\newcommand{\D}{{\mathrm{d}}}
\newcommand{\cE}{{\mathcal{E}}}
\newcommand{\cG}{{\mathcal{G}}}
\newcommand{\cX}{{\mathcal{X}}}
\newcommand{\cY}{{\mathcal{Y}}}
\newcommand{\cB}{{\mathcal{B}}}
\newcommand{\cF}{{\mathcal{F}}}
\newcommand{\cO}{{\mathcal{O}}}
\newcommand{\cS}{{\mathcal{S}}}
\newcommand{\cT}{{\mathcal{T}}}
\newcommand{\given}{\, \big| \,}
\newcommand{\cL}{{\mathscr{L}}}
\newcommand{\cD}{{\mathcal{D}}}
\newcommand{\cA}{{\mathscr{A}}}
\newcommand{\cP}{{\mathscr{P}}}
\newcommand{\cR}{{\mathscr{R}}}
\newcommand{\A}{{\mathcal A}}
\newcommand{\one}{\mathbb{1}}
\newcommand{\mathbbm}[1]{\mathbb{#1}}

\providecommand{\C}{}
\renewcommand{\C}{\mathbb C}
\newcommand{\R}{\mathbb R}
\newcommand{\Q}{\mathbb Q}
\newcommand{\N}{\mathbb N}
\newcommand{\T}{\mathbb T}
\newcommand{\F}{\mathbb F}
\renewcommand{\P}{\mathbb P}
\newcommand{\Z}{\mathbb Z}
\newcommand{\E}{\mathbb {E}}
\newcommand{\Pbb}[1]{\P\left( #1\right)}
\newcommand{\Ebb}[1]{\E\left[ #1\right]}

\newcommand{\lb}{\left (}
\newcommand{\rb}{\right )}
\newcommand{\lbb}{\left [}
\newcommand{\rbb}{\right ]}
\newcommand{\labs}{\left |}
\newcommand{\rabs}{\right |}
\newcommand{\lbrb}[1]{\lb #1 \rb}
\newcommand{\lbbrbb}[1]{\lbb#1\rbb}
\newcommand{\labsrabs}[1]{\labs#1\rabs}
\newcommand{\lbbrb}[1]{\lbb#1\rb}
\newcommand{\lbrbb}[1]{\lb#1\rbb}
\newcommand{\langlerangle}[1]{\langle#1\rangle}
\newcommand{\lbcurly}{\{}
\newcommand{\rbcurly}{\}}
\newcommand{\curly}[1]{\lbcurly#1\rbcurly}
\newcommand{\abs}{\labsrabs}
\newcommand{\ind}[1]{1_{#1}}
\newcommand{\indep}{\perp \!\!\! \perp}
\renewcommand{\d}{\, \mathrm{d}}

\newtheorem{theorem}{Theorem}
\newtheorem{problemCnt}{Lemma}
\theoremstyle{definition}
\newtheorem{problem}[problemCnt]{Problem}
\newtheorem{problembook}[problemCnt]{\faBook \,  Problem}
\newtheorem{problempencil}[problemCnt]{\faPencil* \,  Problem}
\newtheorem{problemstar}[problemCnt]{\faStar\,  Problem}
\newtheorem{problemstarpencil}[problemCnt]{\faStar \, \faPencil* \,   Problem}
\newtheorem{probleminteresting}[problemCnt]{\faSmile[regular] \,   Problem}
\newtheorem{solution}{Solution}

\newcommand{\fwcell}[1]{\makebox[\arraystretch\normalbaselineskip]{$#1$}}
\newenvironment{enumeratei}{\begin{enumerate}[(i)]}{\end{enumerate}}
\newenvironment{enumeratea}{\begin{enumerate}[(a)]}{\end{enumerate}}
\newenvironment{enumeraten}{\begin{enumerate}[\upshape 1.]}{\end{enumerate}}

% Compatibility helpers needed by a few files in the combined archive.
\AtBeginDocument{\noindent{\footnotesize English translation prepared with ChatGPT (OpenAI, GPT-6), 26 October 2026, for translation of the Bulgarian text into English and LaTeX formatting.
The Bulgarian originals remain the reference; translation may contain errors.}\par\smallskip}
\newcommand{\MissingGraphic}[1]{%
  \begin{center}\fbox{\parbox{0.82\linewidth}{\centering Missing source image: \texttt{#1}}}\end{center}}



\setlength{\emergencystretch}{3em}
\allowdisplaybreaks
\begin{document}


\noindent Probability and Statistics, Groups 4/5, 18.12.2021 \hfill Test 1

\hrulefill

\begin{it}
Your grade is 2 plus the number of points earned. Time allowed: 3 hours. Good luck.

Throughout, we work on a probability space $(\Omega, \cF, \P)$.

\end{it}

\begin{problem} (2 points)
According to the data used in this problem, about 2\% of the population has COVID-19. Test manufacturers
report that rapid tests are positive with probability 65\% for infected people and 1\%
for healthy people.

\begin{enumerate}
    \item (0.5 points) A randomly selected person tests positive. What is the probability they are infected?

    $\bullet$
    Let $Inf=\{\text{the person is infected}\}$ and $Pos=\{\text{the test is positive}\}$.
    The required probability is
    \[
    \P(Inf|Pos) = \frac{\P(Pos|Inf)\P(Inf)}{\P(Pos|Inf)\P(Inf) +
    \P(Pos|\overline{Inf})\P(\overline{Inf})}
    = \frac{65\% \times 2 \%}{65\% \times 2 \% + 1 \% \times 98 \%} \approx 57.02 \% \, .
    \]


    \item (0.5 points) In the setting of this problem, a negative test is needed for a green certificate. Suppose a person knows they are infected. How many tests on average are needed until a negative result\footnote{For this problem, ignore the requirement to quarantine automatically after a positive result.}? What is the probability that two tests suffice to obtain the certificate?

    $\bullet$ The number of tests until the first positive result for an infected person is $X\sim Ge(p)$ with $p=35/100$. Thus, on average,
    $100/35\approx2.86$ tests are needed until the first negative result, and the answer to the second part is
    $\P(X \leq 2) = \P(X=1) + \P(X=2) = p + (1-p)p = 0.5775$.

    \item (1 point) Consider two randomly selected people. What is the most likely pair (number infected, number of positive tests)? Calculate the correlation between
    these two counts, using approximations to accuracy 0.01.

    $\bullet$ Let $X$ be the number infected, $Y$ the number of positive tests, $p=99/100$, $q=65/100$
    and $r=2/100$. We calculate
    \begin{align*}
        &\P(X = 0, Y=0) =
        \P(X=0)\P(Y=0|X=0) =
        (1-r)^2 p^2\approx 0.94\\
        &\P(X = 0, Y=1) = (1-r)^2\cdot 2p(1-p)        \approx 0.02\\
        &\P(X = 0, Y=2) \approx 0.00\\
        &\P(X = 1, Y=0) = 2r(1-r)p(1-q)\approx 0.01\\
        &\P(X = 1, Y=1) = 2r(1-r)(pq + (1-p)(1-q))\approx0.03\\
        &\P(X = 1, Y=2) = 2r(1-r)q(1-p)\approx0.00\\
        &\P(X = 2, Y=0) = \P(X=2)\P(Y=0|X=2) = r^2(1-q)^2\approx 0.00\\
        &\P(X = 2, Y=1) = r^2\cdot 2q(1-q)\approx 0.00\\
        &\P(X = 2, Y=2) = r^2q^2\approx 0.00\\
    \end{align*}

Hence the most likely outcome is two healthy people and zero positive tests, and the approximate joint distribution is
    \begin{center}
\begin{tabular}{ |c|c|c|c|c| }
 \hline
$X \backslash Y$ & 0 & 1 & 2& $\P(X = \cdot)$ \\
  \hline
 0 & 0.94& $ 0.02$ &   0.00&0.96\\
 1 &  0.01 &  0.03 &  0.00&0.04\\
 2 & 0.00& 0.00& 0.00&0.00\\
 \hline
 $\P(Y = \cdot)$ &0.95& 0.05 &  0.00& \\
  \hline
  \end{tabular} \, ,
 \end{center}

Successively calculating gives

\begin{align*}
    &\E X \approx 0.04;    &D X = \E X^2 - (\E X)^2 \approx 0.04 - (0.04)^2 \approx 0.04;\\
    &\E Y \approx 0.05;    &D Y =  \E Y^2 - (\E Y)^2 \approx 0.05-0.0025 \approx 0.05;\\
    &\E XY \approx 0.03;
    &Cor(X,Y) = \frac{\E XY - \E X \E Y}{\sqrt{D X \cdot D Y}} \approx
    \frac{0.03 - 0.04\cdot0.05 }{\sqrt{0.04\cdot0.05}} \approx 0.63 \, .
\end{align*}

\end{enumerate}

\end{problem}


\begin{problem} (1.5 points)
A casino offers this game: roll a die 100 times and receive BGN 1 for each six. The entry fee is BGN 20.

\begin{itemize}
    \item (0.5 points) Let
$X_i$ be the winnings from rolls $25i-24$ through $25i$
for $i=1,2,3,4$, and let $X$ be the total winnings. What are the distributions of these
random variables? Would you play?

$\bullet$ $X_i\sim Bin(25,1/6)$ and $X\sim Bin(100,1/6)$. $\E X=100/6<20$, so
the game is not fair in terms of expected profit.

\item (0.25 points) Let $Y=4X_1$: instead of rolling 100 times, roll 25 times and multiply
the winnings by 4. Compare the expectations and variances of $X,Y$, and briefly explain
the result.

$\bullet$
$\E X=\E X_1$, $D X=500/36$, $D(4X_1)=16D X_1=2000/36$. With more
rolls, the deviation from the mean is naturally smaller.

\item (0.75 points)
The casino offers an alternative game: if the sum of 100 rolls exceeds 555, you win BGN 1000. The entry fee is BGN 10.
An advertisement calls this a special Christmas offer and claims the winning probability exceeds $1/100$, giving positive expected profit. Can you confirm or refute this?

$\bullet$ Let $Y_i$ be the value of roll $i$. We can
write the total $Y$ as $Y_1+\dots+Y_{100}$. We calculate
$\E Y_1=7/2$ and $D Y_1=35/12$. Since the $Y_i$ are independent,
$\E Y=100\cdot7/2=350$ and $D Y=100\cdot35/12=875/3$. Thus,
by Chebyshev's inequality,

\[
\P(Y \geq 555) = \P(Y - \E y \geq 205) \leq \P(|Y - \E Y| \geq 205)
\leq \frac{D Y}{205^2} < 0.7\% \, ,
\]
so the winning probability is strictly less than 1\% and the advertisement is false.



\end{itemize}


\end{problem}

\begin{problem} (0.5 points)
Two points are chosen at random on a unit segment. Call its endpoints left and right. What is the probability
that the three resulting segment lengths increase from left to right?
\end{problem}

$\bullet$ Let $x,y$ be the distances of the two points from the left endpoint. If
$y>x$, we need $x<y-x<1-y$, or $2x<y<(x+1)/2$; similarly, if $x>y$, we need $2x-1<y<x/2$.

The final answer follows by elementary geometry or by
\[
2 \int_0^{1/3} \frac{x+1}{2} - 2x \d x = \int_0^{1/3} 1- 3x \d x
= \frac13 - \frac16 = \frac16 \, .
\]

\begin{center}

\end{center}
































































































































































































































































\end{document}
